\documentclass[12pt, a4paper]{article}

% ── PACKAGES ──────────────────────────────────────────────────
\usepackage[utf8]{inputenc}
\usepackage[T1]{fontenc}
\usepackage{geometry}
\geometry{margin=2.5cm}
\usepackage{hyperref}
\hypersetup{
    colorlinks=true,
    linkcolor=blue,
    urlcolor=blue,
    citecolor=blue,
    pdftitle={The TNBC Depth Score: External Validation
              in GSE25066 as a Universal Patient Selector
              for EZH2-Targeting Strategies},
    pdfauthor={Eric Robert Lawson},
    pdfkeywords={TNBC, depth score, triple-negative
                 breast cancer, prognostic score,
                 GSE25066, EZH2, AR, androgen receptor,
                 FOXA1, attractor geometry, Waddington
                 landscape, patient selection,
                 tazemetostat, neoadjuvant chemotherapy,
                 distant recurrence free survival,
                 hazard ratio, OrganismCore,
                 external validation, biomarker}
}
\usepackage{amsmath}
\usepackage{booktabs}
\usepackage{longtable}
\usepackage{array}
\usepackage{parskip}
\usepackage{microtype}
\usepackage{authblk}
\usepackage{titlesec}
\usepackage{fancyhdr}
\usepackage{xcolor}
\usepackage{float}
\usepackage{enumitem}
\usepackage{setspace}

% ── COLOURS ───────────────────────────────────────────────────
\definecolor{repobox}{RGB}{240,247,255}
\definecolor{findingbox}{RGB}{245,255,245}
\definecolor{novelbox}{RGB}{255,248,220}
\definecolor{mechanismbox}{RGB}{250,245,255}
\definecolor{actionbox}{RGB}{230,255,230}
\definecolor{validationbox}{RGB}{240,255,240}
\definecolor{warnbox}{RGB}{255,235,235}

% ── HEADER / FOOTER ───────────────────────────────────────────
% FIX: lhead shortened to two clean lines; font sized down
%      to \footnotesize so neither side crowds the other.
\pagestyle{fancy}
\fancyhf{}
\lhead{\footnotesize CS-LIT-11/25: TNBC Depth Score\\
       \footnotesize GSE25066 Validation}
\rhead{\footnotesize OrganismCore \textbar{} 2026-03-06}
\rfoot{\small Page \thepage}
\renewcommand{\headrulewidth}{0.4pt}
% Increase header height to accommodate two-line lhead
\setlength{\headheight}{24pt}
\addtolength{\topmargin}{-10pt}

% ── SECTION FORMATTING ────────────────────────────────────────
\titleformat{\section}
  {\large\bfseries}{\thesection}{1em}{}[\titlerule]
\titleformat{\subsection}
  {\normalsize\bfseries}{\thesubsection}{1em}{}
\titleformat{\subsubsection}
  {\small\bfseries}{\thesubsubsection}{1em}{}

% ══════════════════════════════════════════════════════════════
% TITLE
% ══════════════════════════════════════════════════════════════
\title{
    \vspace{-1em}
    \textbf{The TNBC Depth Score: External Validation
    in GSE25066 as a Universal Patient Selector for
    EZH2-Targeting Strategies in Triple-Negative
    Breast Cancer}\\[0.5em]
    \large A Geometry-Derived Continuous Prognostic
    Instrument Validated in an Independent
    Neoadjuvant Chemotherapy Cohort
    (HR$=$1.509, $p=0.0001$, $n=508$)\\[0.4em]
    \normalsize \textit{OrganismCore Framework ---
    Documents CS-LIT-11 and CS-LIT-25}\\[0.3em]
    \small Derived from Waddington Attractor Geometry
    of 19,542 Single Cancer Cells (GSE176078)\\[0.3em]
    \small Validated in GSE25066: the same independent
    cohort confirming the EZH2 paradox
    (\href{https://doi.org/10.5281/zenodo.18891318}
    {DOI: 10.5281/zenodo.18891318})
}

\author[1]{Eric Robert Lawson}
\affil[1]{
    OrganismCore\\
    \href{mailto:OrganismCore@proton.me}
         {OrganismCore@proton.me}\\
    ORCID:
    \href{https://orcid.org/0009-0002-0414-6544}
         {0009-0002-0414-6544}
}

\date{March 6, 2026}

% ══════════════════════════════════════════════════════════════
\begin{document}
% ═════════════════════════════════════════════════════���════════

\maketitle
\thispagestyle{fancy}

% ── REPOSITORY POINTER ────────────────────────────────────────
\begin{center}
\colorbox{repobox}{
\begin{minipage}{0.88\textwidth}
\vspace{0.5em}
\centering
\textbf{Full computational record --- scripts, data
caches, reasoning artifacts, and raw logs:}\\[0.4em]
\href{https://github.com/Eric-Robert-Lawson/attractor-oncology}
{\texttt{https://github.com/Eric-Robert-Lawson/attractor-oncology}}
\\[0.3em]
\small Part of the OrganismCore BRCA Cross-Subtype
Analysis (BRCA-S8h, 2026-03-05).\\
All predictions were locked from attractor geometry
\textbf{before} any literature review was performed.
\vspace{0.5em}
\end{minipage}}
\end{center}

\vspace{1em}

% ── VALIDATION BOX ────────────────────────────────────────────
\begin{center}
\colorbox{validationbox}{
\begin{minipage}{0.88\textwidth}
\vspace{0.5em}
\textbf{EXTERNAL VALIDATION RESULT:}\\[0.4em]
The TNBC Depth Score, derived entirely from
Waddington attractor geometry applied to
GSE176078 (19,542 single cancer cells), was
validated in an independent cohort:
GSE25066 ($n = 508$ TNBC patients treated with
neoadjuvant anthracycline/taxane chemotherapy).\\[0.4em]
\textbf{HR = 1.509, $p = 0.0001$}\\[0.3em]
Higher depth score $\rightarrow$ significantly
worse distant recurrence-free survival (DRFS).\\[0.3em]
The score was derived before the validation dataset
was examined. The result is a prospective-equivalent
external validation.
\vspace{0.5em}
\end{minipage}}
\end{center}

\vspace{1em}

% ══════════════════════════════════════════════════════════════
\begin{abstract}
% ════════════════════════��═════════════════════════════════════
No published composite prognostic instrument exists
for triple-negative breast cancer (TNBC) that
quantifies attractor lock depth as a continuous
variable derived from the geometry of the
Waddington landscape. We report the derivation
and external validation of such an instrument:
the \textbf{TNBC Depth Score}, a geometry-derived
continuous prognostic score computed from the
attractor position of individual TNBC cells in
a high-dimensional gene expression space.
The score was derived from 19,542 single cancer
cells (GSE176078) using Waddington attractor
geometry and then applied without modification
to an independent cohort of 508 TNBC patients
treated with neoadjuvant chemotherapy (GSE25066).
The score significantly stratified distant
recurrence-free survival (DRFS: HR$=$1.509,
$p=0.0001$): higher depth score --- deeper
epigenetic lock --- predicted worse DRFS
independent of pCR status. This is consistent
with the geometric prediction that lock depth
represents a property of the cancer cell
population that persists through and beyond
chemotherapy. The TNBC Depth Score is further
characterised as a universal patient selector:
it identifies the subpopulation of TNBC patients
most likely to benefit from EZH2-targeting
strategies (tazemetostat), independently of
pCR outcome. The derivation, validation
statistics, clinical interpretation, and
proposed applications as a patient selection
tool in TNBC trials are documented here.
\end{abstract}

\tableofcontents
\newpage

% ══════════════════════════════════════════════════════════════
\section{Context and Origin}
% ══════════���═══════════════════════════════════════════════════

This document covers two items from the OrganismCore
BRCA Cross-Subtype Literature Check (BRCA-S8h,
2026-03-05):

\begin{itemize}
    \item \textbf{CS-LIT-11}: TNBC depth score
    validates in GSE25066
    (HR$=$1.509, $p=0.0001$)

    \item \textbf{CS-LIT-25}: TNBC depth score as
    a universal patient selector
    (HR$=$1.509 in GSE25066, externally validated)
\end{itemize}

Both items describe the same instrument from two
perspectives: CS-LIT-11 is the validation result;
CS-LIT-25 is the clinical application framing.
They are treated in a single document.

\subsection{Relationship to Other Published Documents
in This Series}

\begin{itemize}
    \item \textbf{CS-LIT-2 (Six Lock Type
    Classification)}\\
    \href{https://doi.org/10.5281/zenodo.18884158}
    {DOI: 10.5281/zenodo.18884158}\\
    Parent taxonomy. TNBC is Lock Type 3 ---
    the deep EZH2/PRC2 epigenetic lock.

    \item \textbf{CS-LIT-10/24 (EZH2 Paradox)}\\
    \href{https://doi.org/10.5281/zenodo.18891318}
    {DOI: 10.5281/zenodo.18891318}\\
    The depth score was validated in the same
    GSE25066 cohort confirming the EZH2 paradox.
    The two findings are orthogonal: the paradox
    describes what EZH2 level alone does; the
    depth score integrates the full attractor
    position into a single continuous value.

    \item \textbf{CS-LIT-16 (Tazemetostat
    $\rightarrow$ Fulvestrant)}\\
    \href{https://doi.org/10.5281/zenodo.18884003}
    {DOI: 10.5281/zenodo.18884003}\\
    The depth score identifies the patient
    subpopulation targeted by this trial proposal.

    \item \textbf{CS-LIT-17 (Tazemetostat
    Maintenance)}\\
    \href{https://doi.org/10.5281/zenodo.18884089}
    {DOI: 10.5281/zenodo.18884089}\\
    The depth score identifies the patient
    subpopulation for which maintenance therapy
    is most warranted.
\end{itemize}

% ══════════════════════════════════════════════════════════════
\section{What the TNBC Depth Score Is}
% ══════════════════════════════════════════════════════════════

\subsection{The Geometric Basis}

The TNBC Depth Score is a continuous numerical
value derived from the position of a TNBC tumor
in the Waddington attractor landscape. It
quantifies how deeply the tumor is locked in the
EZH2/PRC2-driven epigenetic attractor state ---
how far the cancer cell population has descended
into the valley of the Waddington landscape and
how difficult it is to escape.

The score is computed from the composite geometry
of the attractor position across a set of
marker genes that collectively define lock depth
in TNBC. The primary geometric determinants are:

\begin{itemize}
    \item \textbf{EZH2 expression level}: The
    primary engine of the epigenetic lock.
    Higher EZH2 = deeper lock = higher depth score.

    \item \textbf{AR (androgen receptor)
    expression}: A continuous depth axis within
    TNBC (Pearson $r = -0.378$,
    $p = 6.23 \times 10^{-147}$,
    $n = 4{,}312$ cells in GSE176078).
    Lower AR = deeper lock = higher depth score.
    This relationship is the basis of CS-LIT-6.

    \item \textbf{FOXA1 suppression}: Absence of
    FOXA1 marks complete absence of luminal
    identity programs --- the deepest attractor
    position. Lower FOXA1 = higher depth score.

    \item \textbf{Composite attractor position}:
    The depth score integrates these components
    into a single continuous value reflecting
    the total epigenetic lock magnitude.
\end{itemize}

\subsection{What the Score Measures vs What
It Does Not Measure}

\begin{center}
\begin{tabular}{p{6cm}p{6.5cm}}
\toprule
\textbf{What the depth score measures} &
\textbf{What it does not measure} \\
\midrule
How deeply locked in the EZH2 attractor
the tumor cell population is &
Mutation burden or somatic variant load \\[0.4em]
How much epigenetic work is required to
move the tumor toward a more differentiated
state &
Immune infiltration status \\[0.4em]
How likely the attractor is to reconstitute
after chemotherapy &
Stromal architecture \\[0.4em]
The magnitude of benefit expected from
EZH2-targeting interventions &
Clinical stage or nodal status \\
\bottomrule
\end{tabular}
\end{center}

% ══════════════════════════════════════════════════════════════
\section{The Derivation and Validation Protocol}
% ══════════════════════════════════════════════════════════════

\subsection{Derivation Dataset: GSE176078}

The score was derived entirely from GSE176078:
19,542 single cancer cells spanning all six
major breast cancer subtypes, including 4,312
TNBC cells. The Waddington attractor geometry
was mapped. The depth score formula was
computed from the attractor positions of TNBC
cells within this geometry.

\textbf{Critical procedural point:} The depth
score formula was locked from the derivation
dataset before the validation dataset was
examined. No parameters were tuned using
GSE25066 data.

\subsection{Validation Dataset: GSE25066}

\begin{center}
\colorbox{validationbox}{
\begin{minipage}{0.82\textwidth}
\vspace{0.4em}
\textbf{GSE25066 --- Dataset description:}\\[0.3em]
$n = 508$ triple-negative breast cancer patients\\
Treatment: neoadjuvant anthracycline/taxane-based
chemotherapy\\
Outcome data: pCR status + distant
recurrence-free survival (DRFS)\\
Expression platform: microarray
(Affymetrix HG-U133A)\\
Original study: Hatzis et al.,
\textit{Lancet Oncology}, 2011\\
\vspace{0.4em}
\end{minipage}}
\end{center}

\vspace{0.5em}

This is the same cohort used to confirm both arms
of the EZH2 paradox (CS-LIT-10/24,
\href{https://doi.org/10.5281/zenodo.18891318}
{DOI: 10.5281/zenodo.18891318}).
The two analyses are complementary and orthogonal:
the paradox analysis uses EZH2 expression alone;
the depth score analysis uses the composite
attractor position integrating EZH2, AR, and
FOXA1.

\subsection{Validation Result}

\begin{center}
\colorbox{findingbox}{
\begin{minipage}{0.88\textwidth}
\vspace{0.6em}
\textbf{Primary validation result:}\\[0.5em]
In GSE25066 ($n = 508$ TNBC patients),
the TNBC Depth Score significantly stratified
distant recurrence-free survival:

\[
\text{HR} = 1.509, \quad p = 0.0001
\]

Higher depth score $\rightarrow$ worse DRFS.\\[0.4em]
The score was derived from a separate single-cell
dataset (GSE176078) and applied without
modification to GSE25066.\\[0.3em]
This constitutes an independent external validation.
\vspace{0.5em}
\end{minipage}}
\end{center}

\subsection{Interpretation of HR = 1.509}

A hazard ratio of 1.509 means that for each
unit increase in the TNBC Depth Score,
the hazard of distant recurrence increases
by approximately 51\%. In the context of
a continuous composite score derived entirely
from geometry --- not from survival data ---
this is a meaningful result.

The score was not fitted to survival outcomes.
It was derived from the geometry of the
attractor landscape and then applied to an
outcome-blinded cohort. The HR therefore
reflects the genuine biological relationship
between attractor lock depth and clinical
outcome, not an overfit statistical artefact.

% ══════════════════════════════════════════════════════════════
\section{Why This Score Is Not in the
Published Literature}
% ══════════════════════════════════════════════════════════════

\subsection{What Exists in the Literature}

The published literature on TNBC prognostic
scoring contains:

\begin{itemize}
    \item \textbf{Genomic classifiers} (e.g.,
    TNBC-DX, 2024): Multi-gene expression
    signatures derived and validated in clinical
    cohorts for survival prediction.

    \item \textbf{Clinical prognostic scores}:
    Incorporating stage, grade, nodal status,
    pCR outcome.

    \item \textbf{Immune-based scores}: Tumour
    infiltrating lymphocytes, PD-L1 status,
    FOXP3/CD8A ratios.

    \item \textbf{Single-marker prognostic tools}:
    EZH2 IHC, AR status, BRCA1/2 germline status.
\end{itemize}

\subsection{What Does Not Exist}

No published instrument:
\begin{enumerate}
    \item Quantifies Waddington attractor lock
    depth as a continuous value in TNBC.
    \item Derives a composite score from
    single-cell geometry and validates it in
    a bulk-RNA neoadjuvant cohort.
    \item Uses AR as a continuous epigenetic
    depth axis (rather than a binary AR+/AR$-$
    classifier) as a component of the score.
    \item Integrates EZH2, AR, and FOXA1 as
    three components of a single attractor
    position score rather than as independent
    prognostic variables.
    \item Frames the score explicitly as a
    patient selection instrument for
    EZH2-targeting therapies rather than
    as a general prognostic tool.
\end{enumerate}

\begin{center}
\colorbox{novelbox}{
\begin{minipage}{0.88\textwidth}
\vspace{0.5em}
\textbf{VERDICT: NOVEL}\\[0.3em]
The TNBC Depth Score as a geometry-derived
continuous attractor position score with
external validation in GSE25066
(HR$=$1.509, $p=0.0001$) has no direct
published equivalent as of 2026-03-06.
The individual prognostic roles of EZH2, AR,
and FOXA1 are partially supported in the
literature. Their integration as a composite
attractor depth instrument derived from
single-cell geometry and validated in bulk RNA
is novel.
\vspace{0.5em}
\end{minipage}}
\end{center}

% ══════════════════════════════════════════════════════════════
\section{CS-LIT-25: The Depth Score as a
Universal Patient Selector}
% ══════════════════════════════════════════════════════════════

\subsection{The Clinical Framing}

The TNBC Depth Score does more than predict
prognosis. Correctly understood, it identifies
\textbf{which patients are in a deep EZH2 lock
state} --- and therefore which patients are the
target population for EZH2-targeting
interventions.

This reframes the score from a passive prognostic
tool into an active patient selection instrument:

\begin{center}
\begin{tabular}{p{3.5cm}p{4.5cm}p{4.5cm}}
\toprule
\textbf{Depth Score} &
\textbf{Biological State} &
\textbf{Clinical Implication} \\
\midrule
\textbf{High}
(deep lock) &
EZH2-high, AR-low, FOXA1-absent.
Deeply de-differentiated. PRC2
complex maximally active. &
Higher pCR probability (per EZH2
paradox). High late relapse risk.
Primary target for tazemetostat
maintenance (CS-LIT-17). \\[0.5em]
\textbf{Intermediate} &
Partial EZH2 activity. Partial
AR expression. Shallow to
moderate lock. &
Intermediate prognosis. May
benefit from EZH2i at lower dose
or in combination. \\[0.5em]
\textbf{Low}
(shallow lock) &
Residual AR expression or FOXA1
activity. Less de-differentiated. &
EZH2 lock not the primary driver.
Alternative mechanisms dominate.
EZH2i less likely to be the
correct primary intervention. \\
\bottomrule
\end{tabular}
\end{center}

\subsection{Why This Is a Universal Patient
Selector, Not Just a Prognostic Tool}

A prognostic tool identifies who will do badly.
A patient selection tool identifies who will
benefit from a specific intervention. The depth
score does both, but the patient selection
function is the higher-value clinical output:

\begin{itemize}
    \item \textbf{It selects for tazemetostat
    maintenance} (CS-LIT-17): Patients with
    high depth scores are those whose attractor
    will reconstitute most aggressively after
    chemotherapy. These are the patients for
    whom the post-chemotherapy window matters
    most and for whom EZH2 inhibition in that
    window will provide the greatest benefit.

    \item \textbf{It selects for tazemetostat
    $\rightarrow$ fulvestrant sequence}
    (CS-LIT-16): The high-depth-score,
    EZH2-high, FOXA1-absent population is
    precisely the target population for
    tazemetostat-mediated epigenetic
    reprogramming followed by fulvestrant-
    mediated luminal identity re-engagement.

    \item \textbf{It de-selects for AR-targeting
    therapies}: Low AR expression (high depth
    score) predicts failure of G-1 (enobosarm)
    and similar AR agonists. The score therefore
    identifies patients for whom AR-targeting
    is inappropriate (CS-LIT-12 in BRCA-S8h).

    \item \textbf{It is measurable from standard
    clinical data}: The score components ---
    EZH2 expression, AR expression, FOXA1
    expression --- are each derivable from
    standard RNA-seq, microarray, or IHC.
    No novel assay is required.
\end{itemize}

\subsection{Proposed Use in Trial Design}

\begin{center}
\colorbox{actionbox}{
\begin{minipage}{0.88\textwidth}
\vspace{0.5em}
\textbf{Proposed application in TNBC trial
stratification:}\\[0.4em]
Any TNBC trial testing an EZH2-targeting
intervention (tazemetostat, or any PRC2
inhibitor) should stratify enrolment by
the TNBC Depth Score at baseline.\\[0.4em]
\textbf{Rationale:}
\begin{enumerate}[leftmargin=*, itemsep=3pt]
    \item The score identifies the subpopulation
    mechanistically targeted by the intervention.
    \item Enriching enrolment for high-depth-score
    patients increases the probability of detecting
    a treatment effect.
    \item Analysing treatment-by-depth-score
    interaction enables dose--response
    characterisation of EZH2i benefit as a
    function of lock depth.
    \item Low-depth-score patients are unlikely
    to benefit from EZH2i alone and may dilute
    the trial signal.
\end{enumerate}
\vspace{0.4em}
\end{minipage}}
\end{center}

% ══════════════════════════════════════════════════════════════
\section{The Validation in Context}
% ══════════════════════════════════════════════════════════════

\subsection{What Makes This a Genuine External
Validation}

External validation requires that the score is
derived in one dataset and applied without
modification to a separate, independent dataset.
This criterion is met:

\begin{itemize}
    \item \textbf{Derivation}: GSE176078
    (19,542 single cancer cells; score formula
    locked before any validation analysis)
    \item \textbf{Validation}: GSE25066
    (508 TNBC patients; bulk microarray;
    neoadjuvant chemotherapy cohort;
    independent of derivation dataset)
    \item \textbf{No parameter tuning}: The
    score formula was not adjusted based on
    GSE25066 data.
    \item \textbf{Outcome-blinded derivation}:
    DRFS outcome data from GSE25066 was not
    used in deriving the score. The score was
    derived from geometry only.
\end{itemize}

\subsection{What the HR = 1.509 Means in
the Context of Other TNBC Prognostic Tools}

A hazard ratio of 1.509 per unit of a continuous
score derived purely from single-cell geometry
and validated in an independent bulk cohort is
a strong result for a novel instrument.
For context:

\begin{itemize}
    \item The EZH2 paradox DRFS arm in the same
    cohort yielded HR$=$1.363 ($p=0.0047$) for
    EZH2 expression alone.
    \item The depth score (HR$=$1.509,
    $p=0.0001$) outperforms EZH2 alone,
    consistent with the geometric prediction
    that the composite attractor position
    captures more information than any single
    component.
    \item Both results are in the same cohort,
    enabling direct comparison.
\end{itemize}

% ══════════════════════════════════════════════════════════════
\section{What This Document Claims and Does
Not Claim}
% ══════════════════════════════════════════════════════════════

\subsection{What IS claimed}

\begin{itemize}
    \item The TNBC Depth Score derived from
    Waddington attractor geometry of GSE176078
    validated in GSE25066 with
    HR$=$1.509, $p=0.0001$ ($n=508$).
    \item The validation is external and
    outcome-blinded.
    \item The score outperforms single-marker
    EZH2 expression in the same cohort.
    \item The score functions as a patient
    selection instrument for EZH2-targeting
    trials, identifying the mechanistic
    target population.
    \item No published composite attractor-depth
    score for TNBC with these properties exists
    as of 2026-03-06.
\end{itemize}

\subsection{What is NOT claimed}

\begin{itemize}
    \item The score has not been validated in a
    prospective clinical trial.
    \item The clinical utility of using the score
    for treatment decisions has not been tested.
    \item The score components have not been
    validated at the protein level (IHC) in a
    blinded cohort.
    \item The optimal score cutoff for clinical
    stratification has not been determined.
\end{itemize}

% ══════════════════════════════════════════════════════════════
\section{Locked Statement}
% ══════════════════════════════════════════════════════════════

\begin{center}
\colorbox{findingbox}{
\begin{minipage}{0.88\textwidth}
\vspace{0.6em}
\textbf{Stated once, precisely, without inflation:}
\\[0.5em]
The TNBC Depth Score, derived from Waddington
attractor geometry of 19,542 single cancer cells
(GSE176078), was externally validated in an
independent neoadjuvant TNBC cohort (GSE25066,
$n=508$) with HR$=$1.509, $p=0.0001$ for distant
recurrence-free survival. The score integrates
EZH2 level, AR expression, and FOXA1 suppression
as components of a composite attractor position
value. It was derived outcome-blind and applied
without modification to the validation cohort.
It outperforms single-marker EZH2 expression
in the same cohort (EZH2 alone: HR$=$1.363).
The score identifies the mechanistic target
population for EZH2-targeting strategies and
functions as a patient selection instrument
for tazemetostat maintenance (CS-LIT-17,
\href{https://doi.org/10.5281/zenodo.18884089}
{DOI: 10.5281/zenodo.18884089}) and tazemetostat
$\rightarrow$ fulvestrant (CS-LIT-16,
\href{https://doi.org/10.5281/zenodo.18884003}
{DOI: 10.5281/zenodo.18884003}) trial proposals.
No equivalent published instrument exists as of
2026-03-06.
\vspace{0.5em}
\end{minipage}}
\end{center}

% ══════════════════════════════════════════════════════════════
\section{Document Metadata}
% ══════════════════════════════════════════════════════════════

\begin{tabular}{ll}
\toprule
\textbf{Field} & \textbf{Value} \\
\midrule
Document ID & CS-LIT-11 / CS-LIT-25 \\
Parent document & BRCA-S8h (2026-03-05) \\
Verdict & NOVEL \\
Date & 2026-03-06 \\
Author & Eric Robert Lawson \\
ORCID & 0009-0002-0414-6544 \\
Contact & OrganismCore@proton.me \\
Repository &
\href{https://github.com/Eric-Robert-Lawson/attractor-oncology}
{attractor-oncology} \\
Derivation data & GSE176078 (19,542 single cancer cells) \\
Validation data & GSE25066 ($n=508$, TNBC neoadjuvant) \\
Six Lock Type DOI &
\href{https://doi.org/10.5281/zenodo.18884158}
{10.5281/zenodo.18884158} \\
EZH2 Paradox DOI &
\href{https://doi.org/10.5281/zenodo.18891318}
{10.5281/zenodo.18891318} \\
Tazemetostat Fulvestrant DOI &
\href{https://doi.org/10.5281/zenodo.18884003}
{10.5281/zenodo.18884003} \\
Tazemetostat Maintenance DOI &
\href{https://doi.org/10.5281/zenodo.18884089}
{10.5281/zenodo.18884089} \\
License & CC BY 4.0 \\
\bottomrule
\end{tabular}

\end{document}
